\section{Special-function context and the contribution boundary}
\label{sec:literature}

Polynomial heat evolution and its finite operational expansion are
classical. Hile and Stanoyevitch's framework
\cite[Theorem 1 and (2.18)]{HileStanoyevitch2001} includes polynomial
initial data for constant-coefficient evolution equations.
Batahan and Shehata~\cite[Theorem 3.1]{BatahanShehata2014} give a
Hermite--Chebyshev operational representation with a fixed negative
heat parameter and the second-kind polynomial. Related two-variable
formulas appear in Khan~\cite{Khan2019}. These formulas are cited as
historical context, not as inputs to the positive-time estimates above;
the finite Gaussian representation and the two normalizations used here
are verified directly.

The standard trigonometric definitions of the two Chebyshev families
are recorded in \cite[\S18.5]{DLMF}. At zero displacement there is also
an exact terminating identity:
\begin{equation}\label{eq:terminating-hypergeometric}
 Q_{m,A}(s,0)={}_4F_2\left(\begin{matrix}
 -m/2,(1-m)/2,m/2,(m+1)/2\\1/4,3/4
 \end{matrix};\frac{s}{A^2}\right).
\end{equation}
Here the notation means the finite sum through $j=\lfloor m/2\rfloor$.
To verify it, let $D_k=T_m^{(k)}(1)$ and differentiate the Chebyshev
ODE at one. The recurrence
$D_{k+1}/D_k=(m^2-k^2)/(2k+1)$, $D_0=1$, shows that the ratio of
successive terms $D_{2j}(s/A^2)^j/j!$ is exactly the hypergeometric
term ratio in \eqref{eq:terminating-hypergeometric}.
The constant terms also agree. This is an algebraic representation,
not an additional asymptotic theorem.

Its leading coefficient model is
\[
 {}_0F_2(;1/4,3/4;s\mu^6/16).
\]
The classical exponential factor $\exp(3z^{1/3})$ of this function
explains $F_s\mu^2$; see \cite[\S16.11]{DLMF}.
It does not imply that the two functions have ratio tending to one.
Indeed the exact $j$th coefficient relative to that model contains
\[
 \prod_{k=0}^{2j-1}(1-k^2/m^2).
\]
For $j\asymp\mu^2$ in the original regime, $j/m\to0$, and the
logarithmic correction has order $j^3/m^2\asymp\mu^3/A$.
This follows by summing
$-2x\le\log(1-x)\le-x$ once $(2j/m)^2\le1/2$.
That correction is $o(\mu^2)$ but need not be $o(1)$.

The large-variable expansions in DLMF \S16.11(ii) explicitly exclude
nonpositive integer upper parameters and treat $p\le q+1$.
They therefore do not directly apply to this terminating $p=q+2$
family. The large-parameter formulas in \S16.11(iii) use other
parameter paths. The uniform Gauss hypergeometric expansions of
Paris~\cite{Paris2018,Paris2021} concern ${}_2F_1$ with specified
large-parameter paths and transition regions. They do not state the
terminating four-upper-parameter heat-edge estimate and its bounded
spatial-displacement quotient used here. This comparison does not
rule out other covering results.

Large-degree polynomial heat flow also appears in work on zeros and
random matrices, for example Hall and Ho~\cite{HallHo2025}.
A macroscopic zero distribution differs from the function values and
relative amplitudes at shrinking endpoint distances considered here.
The bounded primary-source search found no theorem with the same
joint regime, explicit uniform errors, and displacement quantifiers.
It is not an exhaustive search of all terminating hypergeometric
literature and establishes no first-result or priority claim.

The claimed addition is the quantitative statement with these
hypotheses, its checked second-family instance, and the fixed-data
full-interval realization below. The Laplace principle, Gaussian
tilting, and polynomial heat operation remain established methods.
